\begin{lemma}\label{lem:16.2}
	Let $r\in\N$ and $n_1,\dotsc,n_r\in\N_0$. Let $\gamma$ be a partition and for each $i\in\{1,2,\dotsc,r\}$, let $\psi_i$ be a virtual character of $S_{n_i}$. Then
	\[ (\psi_1\boxtimes\cdots\boxtimes\psi_r)/\chi^\gamma = \sum_{\gamma_1,\dotsc,\gamma_r} c^\gamma_{\gamma_1,\dotsc,\gamma_r}\cdot (\psi_1/\chi^{\gamma_1})\boxtimes\cdots\boxtimes(\psi_r/\chi^{\gamma_r}), \]
	summed over all sequences of partitions $\gamma_1,\dotsc,\gamma_r$ such that $|\gamma_1|+\cdots+|\gamma_r|=|\gamma|$.
\end{lemma}

\begin{proof}
	The case $r=1$ is trivial. For ease of notation we prove the statement for $r=2$; the case of general $r$ follows by an analogous argument. Define $n\coloneqq n_1+n_2$, $k\coloneqq |\gamma|$ and let $\delta\vdash n-k$. Then
	\begin{align*}
		\langle (\psi_1\boxtimes\psi_2)/\chi^\gamma,\chi^\delta\rangle &= \langle (\psi_1\times\psi_2)\up_{S_{n_1}\times S_{n_2}}^{S_n}, (\chi^\delta\times\chi^\gamma)\up_{S_{n-k}\times S_k}^{S_n}\rangle\\
		&= \langle (\psi_1\times\psi_2)\up_{S_{n_1}\times S_{n_2}}^{S_n}\down_{S_{n-k}\times S_k}, \chi^\delta\times\chi^\gamma\rangle,
	\end{align*}
	and applying Mackey's Theorem,
	\footnotesize
	\begin{align*}
		&= \sum_{\substack{0\le t_1,t_2\le k\\ t_1+t_2=k}} \langle (\psi_1\times\psi_2)\down^{S_{n_1}\times S_{n_2}}_{S_{n_1-t_1}\times S_{t_1}\times S_{n_2-t_2}\times S_{t_2}}\up^{S_{n-k}\times S_{k}}, \chi^\delta\times\chi^\gamma\rangle\\
		&= \sum_{t_1+t_2=k} \sum_{\substack{i=1,2\\\gamma_i\vdash t_i\\\delta_i\vdash n_i-t_i}} c^\gamma_{\gamma_1,\gamma_2}\cdot c^\delta_{\delta_1,\delta_2}\cdot \langle \psi_1\down_{S_{n_1-t_1}\times S_{t_1}} \times \psi_2\down_{S_{n_2-t_2}\times S_{t_2}}, \chi^{\delta_1}\times\chi^{\gamma_1}\times\chi^{\delta_2}\times\chi^{\gamma_2}\rangle\\
		&= \sum_{t_1+t_2=k} \sum_{\substack{i\\\gamma_i\vdash t_i\\\delta_i\vdash n_i-t_i}} c^\gamma_{\gamma_1,\gamma_2}\cdot c^\delta_{\delta_1,\delta_2}\cdot \langle \psi_1/\chi^{\gamma_1},\chi^{\delta_1}\rangle\cdot\langle \psi_2/\chi^{\gamma_2},\chi^{\delta_2}\rangle\\
		&= \sum_{t_1+t_2=k} \sum_{\substack{i\\\gamma_i\vdash t_i}} c^\gamma_{\gamma_1,\gamma_2}\cdot \langle (\psi_1/\chi^{\gamma_1})\boxtimes(\psi_2/\chi^{\gamma_2}), \chi^\delta\rangle,
	\end{align*}
	\normalsize
	as claimed.
\end{proof}

\begin{corollary}\label{cor:16.2}
	Let $m,n\in\N$ and $\lambda=(\lambda_1,\dotsc,\lambda_r)\vdash n$. Let $k\in\{0,1,\dotsc,n-1\}$ and $\beta\vdash k$. Then
	\begin{enumerate}[label=\textup{(\roman*)}]
		\item \label{591} $\big(\operatorname*{\boxtimes}_{i=1}^r \rho^{(\lambda_i)}_m \big)/\chi^{(1^{n-k})} = \sum_{\underline{t}}  \operatorname*{\boxtimes}_{i=1}^r \big( \rho^{(\lambda_i)}_m/\chi^{(1^{\lambda_i-t_i})} \big)$, and
		
		\item \label{592} $\big( \operatorname*{\boxtimes}_{i=1}^r \chi^{(1^{\lambda_i})} \big)/\chi^{\beta'} = \sum_{\underline{t}} c^\beta_{(t_1),\dotsc,(t_r)}\cdot \operatorname*{\boxtimes}_{i=1}^r \chi^{(1^{\lambda_i-t_i})}$,
	\end{enumerate}
	summed over all compositions $\underline{t}=(t_1,\dotsc,t_r)$ of $k$ into $r$ parts. That is, $t_i\in\N_0$ for all $i$ and $t_1+\cdots+t_r=k$ \textup{(}and we may further assume $t_i\le\lambda_i$ for all $i$\textup{)}.
\end{corollary}

\begin{proof}
	\ref{591} Applying Lemma~\ref{lem:16.2} with $\gamma=(1^{n-k})$, observe that $c^\gamma_{\gamma_1,\dotsc,\gamma_r}\in\{0,1\}$ and is non-zero only if each $\gamma_i=(1^{s_i})$ for some $s_i\in\N_0$. By Lemma~\ref{lem:tall-pleth}, $\rho^{(\lambda_i)}_m$ only has irreducible constituents $\chi^\mu$ where $l(\mu)\le \lambda_i$, so we may further assume that $s_i\le\lambda_i$. Writing $t_i=\lambda_i-s_i$ gives the result.
	
	\noindent \ref{592} Applying Lemma~\ref{lem:16.2} with $\psi_i=\chi^{(1^{\lambda_i})}$, observe that $\psi_i/\chi^{\gamma_i}\ne 0$ only if $\gamma_i=(1^{t_i})$ for some $0\le t_i\le \lambda_i$. Moreover, $c^{\beta'}_{(1^{t_1}),\dotsc,(1^{t_r})}=c^\beta_{(t_1),\dotsc,(t_r)}\in\{0,1\}$.
\end{proof}

\begin{lemma}\label{lem:16.3}
	Let $n\in\N$ and $k\in\{0,1,\dotsc,n\}$. Let $\nu$ be a partition with $l(\nu)=n$ and set $\hat{\nu}\coloneqq \nu-(1^n)$. Let $\delta\vdash|\nu|-k$ with $l(\delta)\le n$. Then
	\[ \langle\chi^\delta,\chi^{\nu/(1^k)} \rangle = \langle \chi^{\delta/(1^{n-k})},\chi^{\hat{\nu}}\rangle. \]
\end{lemma}

\begin{proof}
	The determinantal form of skew characters of symmetric groups (see e.g.~\cite[2.3.13]{JK}) gives $\chi^{\alpha/\beta} = \det\big(\chi^{(\alpha_i-i-\beta_j+j)}\big)$ whenever $\alpha$ and $\beta$ are partitions, where the multiplication of characters in expanding the determinant is given by the operation $\boxtimes$. Applying this to $\alpha=\nu'$ and $\beta=\emptyset$ and expanding the determinant with respect to the first row gives
	\[ \chi^{\nu'} = \sum_{j\ge 1} (-1)^{j-1}\cdot \chi^{(n+j-1)}\boxtimes \chi^{\hat{\nu}'/(1^{j-1})} = \sum_{j\ge 0} (-1)^j\cdot\chi^{(n+j)}\boxtimes \chi^{\hat{\nu}'/(1^j)}. \]
	Multiplying both sides by the sign representation then gives
	\[ \chi^\nu = \sum_{j\ge 0} (-1)^j\cdot \chi^{(1^{n+j})}\boxtimes \chi^{\hat{\nu}/(j)}. \]
	Then
	\begin{align*}
		\langle\chi^\delta,\chi^{\nu/(1^k)}\rangle &= \langle \chi^\delta\boxtimes\chi^{(1^k)},\chi^\nu\rangle \\
		&= \sum_{j\ge 0}(-1)^j\cdot\langle(\chi^\delta\boxtimes\chi^{(1^k)})/\chi^{(1^{n+j})}, \chi^{\hat{\nu}/(j)} \rangle\\
		&\overset{\normalfont\tiny \text{(Lemma~\ref{lem:16.2})}}{=} \sum_{j=0}^k (-1)^j\cdot \left\langle \sum_{s=0}^{k-j}(\chi^\delta/\chi^{(1^{n-s})})\boxtimes(\chi^{(1^k)}/\chi^{(1^{s+j})}), \chi^{\hat{\nu}/(j)}\right\rangle\\
		&=\sum_{j=0}^k \sum_{s=0}^{k-j}(-1)^j \langle (\chi^\delta/\chi^{(1^{n-s})})\boxtimes \chi^{(1^{k-j-s})}, \chi^{\hat{\nu}/(j)}\rangle\\
		&= \sum_{s=0}^k \left\langle (\chi^\delta/\chi^{(1^{n-s})})\boxtimes \left(\sum_{j=0}^{k-s}(-1)^j\cdot\chi^{(1^{k-s-j})}\boxtimes\chi^{(j)}\right), \chi^{\hat{\nu}}\right\rangle\\
		&=\langle \chi^\delta/\chi^{(1^{n-k})}, \chi^{\hat{\nu}}\rangle,
	\end{align*}
	where the final equality follows since $\chi^{(1^{k-s-j})}\boxtimes\chi^{(j)}=\chi^{(j+1,1^{k-s-j-1})}+\chi^{(j,1^{k-s-j})}$, and so $\sum_{j=0}^{k-s}(-1)^j\cdot\chi^{(1^{k-s-j})}\boxtimes\chi^{(j)}$ equals zero if $k\ne s$, and equals $\chi^\emptyset$ if $k=s$.
\end{proof}
